% GENERATED FILE. Edit canonical lessons, then run npm run generate:books.
% canonical-source-hash: fnv1a64:78646ab52a40dc61
% canonical-lessons: BN-C544-totpor, BN-C544-obokash, BN-C544-ghoniye-asa, BN-C544-hurohuri, BN-C544-jhotpot

\chapter{Time Pressure}
\label{ch:bn-time-pressure}
\hlchaptermodality{544}

\begin{chapteropening}
\textbf{By the end of this chapter:} \emph{I can talk about rushing, breaks and dates drawing near.}

Complete the last lesson of chapter 544: I can talk about rushing, breaks and dates drawing near.
\end{chapteropening}

\section[\texorpdfstring{\bn{তৎপর} (tôtpôr)}{তৎপর (tôtpôr)}]{\bn{তৎপর} (tôtpôr) — prompt; quick to act}

\label{lesson:BN-C544-totpor}



\begin{quote}
Before the new one: say the Bengali for to secure, then the Bengali for a desired goal.
\end{quote}

\subsection*{You'll want to know: \bn{তৎপর}}
\textbf{\bn{তৎপর}} — \emph{tôtpôr} — \textquotedblleft{}prompt; quick to act\textquotedblright{}.

\begin{quote}
\bn{যদি} \bn{দমকল} \bn{তৎপর} \bn{না} \bn{হত}, \bn{তাহলে} \bn{আগুন} \bn{ছড়াত।}
\end{quote}

\begin{quote}
\emph{Jôdi dômkôl tôtpôr nā hôtô, tāhôle āgun chhôṛātô.}
\end{quote}

\begin{quote}
If the fire brigade hadn't been quick, the fire would have spread.
\end{quote}

\begin{cousinweb}
\textbf{\bn{তৎপর}} (\emph{tôtpôr}) is a tatsama, Sanskrit \emph{tatpara}, \textquotedblleft{}following on that, intent on it\textquotedblright{} (Monier-Williams).
\end{cousinweb}

\begin{grammarlens}[title={What you've built}]
One more word for saying what happens if.
\end{grammarlens}

\subsection*{Your turn}
\begin{itemize}
\raggedright
  \item \emph{Say it:} \emph{tôtpôr}
  \item \emph{Say it:} \emph{tôtpôr}, once more
  \item \emph{Say it:} \emph{Jôdi dômkôl tôtpôr nā hôtô, tāhôle āgun chhôṛātô.} — If the fire brigade hadn't been quick, the fire would have spread.
  \item \emph{Recall:} say the Bengali for to secure, then the Bengali for a desired goal, then say \emph{tôtpôr} again
\end{itemize}

\begin{tcolorbox}[breakable,colback=teal!4,colframe=teal!35!black,title={Before you move on}]
What does \bn{তৎপর} mean? (\textquotedblleft{}Prompt; quick to act\textquotedblright{}.) Say it once more.
\end{tcolorbox}
\section[\texorpdfstring{\bn{অবকাশ} (ôbôkāsh)}{অবকাশ (ôbôkāsh)}]{\bn{অবকাশ} (ôbôkāsh) — a respite; a break}

\label{lesson:BN-C544-obokash}



\begin{quote}
Before the new one: say the Bengali for a desired goal, then the Bengali for prompt.
\end{quote}

\subsection*{You'll want to know: \bn{অবকাশ}}
\textbf{\bn{অবকাশ}} — \emph{ôbôkāsh} — \textquotedblleft{}a respite; a break\textquotedblright{}.

\begin{quote}
\bn{দুপুরে} \bn{অবকাশ} \bn{পেলে} \bn{আমি} \bn{একটু} \bn{ঘুমাই।}
\end{quote}

\begin{quote}
\emph{Dupure ôbôkāsh pele āmi ekṭu ghumāi.}
\end{quote}

\begin{quote}
If I get a break at midday, I take a short nap.
\end{quote}

\begin{cousinweb}
\textbf{\bn{অবকাশ}} (\emph{ôbôkāsh}) is a tatsama, Sanskrit \emph{avakāśa}, \textquotedblleft{}room, space\textquotedblright{}, which Monier-Williams also gives for a stretch of time in between: a gap in the day.
\end{cousinweb}

\begin{grammarlens}[title={What you've built}]
One more word for saying what happens if.
\end{grammarlens}

\subsection*{Your turn}
\begin{itemize}
\raggedright
  \item \emph{Say it:} \emph{ôbôkāsh}
  \item \emph{Say it:} \emph{ôbôkāsh}, once more
  \item \emph{Say it:} \emph{Dupure ôbôkāsh pele āmi ekṭu ghumāi.} — If I get a break at midday, I take a short nap.
  \item \emph{Recall:} say the Bengali for a desired goal, then the Bengali for prompt, then say \emph{ôbôkāsh} again
\end{itemize}

\begin{tcolorbox}[breakable,colback=teal!4,colframe=teal!35!black,title={Before you move on}]
What does \bn{অবকাশ} mean? (\textquotedblleft{}A respite; a break\textquotedblright{}.) Say it once more.
\end{tcolorbox}
\section[\texorpdfstring{\bn{ঘনিয়ে} \bn{আসা} (ghôniye āsā)}{ঘনিয়ে আসা (ghôniye āsā)}]{\bn{ঘনিয়ে} \bn{আসা} (ghôniye āsā) — to draw near (of a date or danger)}

\label{lesson:BN-C544-ghoniye-asa}



\begin{quote}
Before the new one: say the Bengali for prompt, then the Bengali for a respite.
\end{quote}

\subsection*{You'll want to know: \bn{ঘনিয়ে} \bn{আসা}}
\textbf{\bn{ঘনিয়ে} \bn{আসা}} — \emph{ghôniye āsā} — \textquotedblleft{}to draw near (of a date or danger)\textquotedblright{}.

\begin{quote}
\bn{বিয়ের} \bn{দিন} \bn{ঘনিয়ে} \bn{এল।}
\end{quote}

\begin{quote}
\emph{Biyer din ghôniye elô.}
\end{quote}

\begin{quote}
The wedding day drew near.
\end{quote}

Past (he/she …): \bn{সেই} \bn{দিন} \bn{ঘনিয়ে} \bn{এল}, \emph{sei din ghôniye elô}.

\begin{cousinweb}
\textbf{\bn{ঘনিয়ে} \bn{আসা}} (\emph{ghôniye āsā}) joins \bn{ঘনিয়ে}, the -\bn{এ} form of a verb built on \bn{ঘন}, \textquotedblleft{}dense, close\textquotedblright{}, to \bn{আসা}, \textquotedblleft{}to come\textquotedblright{}: coming closer and closer. CDIAL lists Bengali \emph{ghana}, \textquotedblleft{}thick, dense\textquotedblright{}, under Sanskrit \emph{ghana}, and notes a Bengali verb from it meaning to come close (CDIAL 4424).
\end{cousinweb}

\begin{grammarlens}[title={What you've built}]
One more word for saying what happens if.
\end{grammarlens}

\subsection*{Your turn}
\begin{itemize}
\raggedright
  \item \emph{Say it:} \emph{ghôniye āsā}
  \item \emph{Say it:} \emph{ghôniye āsā}, once more
  \item \emph{Say it:} \emph{Biyer din ghôniye elô.} — The wedding day drew near.
  \item \emph{Recall:} say the Bengali for prompt, then the Bengali for a respite, then say \emph{ghôniye āsā} again
\end{itemize}

\begin{tcolorbox}[breakable,colback=teal!4,colframe=teal!35!black,title={Before you move on}]
What does \bn{ঘনিয়ে} \bn{আসা} mean? (\textquotedblleft{}To draw near (of a date or danger)\textquotedblright{}.) Say it once more.
\end{tcolorbox}
\section[\texorpdfstring{\bn{হুড়োহুড়ি} (huṛohuṛi)}{হুড়োহুড়ি (huṛohuṛi)}]{\bn{হুড়োহুড়ি} (huṛohuṛi) — a scramble; a rush}

\label{lesson:BN-C544-hurohuri}



\begin{quote}
Before the new one: say the Bengali for a respite, then the Bengali for to draw near.
\end{quote}

\subsection*{You'll want to know: \bn{হুড়োহুড়ি}}
\textbf{\bn{হুড়োহুড়ি}} — \emph{huṛohuṛi} — \textquotedblleft{}a scramble; a rush\textquotedblright{}.

\begin{quote}
\bn{স্টেশনে} \bn{লোকজনের} \bn{হুড়োহুড়ি} \bn{ছিল।}
\end{quote}

\begin{quote}
\emph{Sṭeshône lokjôner huṛohuṛi chhilô.}
\end{quote}

\begin{quote}
There was a scramble of people at the station.
\end{quote}

\begin{cousinweb}
\textbf{\bn{হুড়োহুড়ি}} (\emph{huṛohuṛi}) is an expressive doubled word for people pushing and rushing. It perhaps belongs with \bn{হুড়}, \textquotedblleft{}a crowd\textquotedblright{}, which CDIAL lists as Bengali \emph{huṛ} under \emph{huḍati}, \textquotedblleft{}collects\textquotedblright{} (CDIAL 14136); the link is not settled.
\end{cousinweb}

\begin{grammarlens}[title={What you've built}]
One more word for saying what happens if.
\end{grammarlens}

\subsection*{Your turn}
\begin{itemize}
\raggedright
  \item \emph{Say it:} \emph{huṛohuṛi}
  \item \emph{Say it:} \emph{huṛohuṛi}, once more
  \item \emph{Say it:} \emph{Sṭeshône lokjôner huṛohuṛi chhilô.} — There was a scramble of people at the station.
  \item \emph{Recall:} say the Bengali for a respite, then the Bengali for to draw near, then say \emph{huṛohuṛi} again
\end{itemize}

\begin{tcolorbox}[breakable,colback=teal!4,colframe=teal!35!black,title={Before you move on}]
What does \bn{হুড়োহুড়ি} mean? (\textquotedblleft{}A scramble; a rush\textquotedblright{}.) Say it once more.
\end{tcolorbox}
\section[\texorpdfstring{\bn{ঝটপট} (jhôṭpôṭ)}{ঝটপট (jhôṭpôṭ)}]{\bn{ঝটপট} (jhôṭpôṭ) — in a jiffy; quickly}

\label{lesson:BN-C544-jhotpot}



\begin{quote}
Before the new one: say the Bengali for to draw near, then the Bengali for a scramble.
\end{quote}

\subsection*{You'll want to know: \bn{ঝটপট}}
\textbf{\bn{ঝটপট}} — \emph{jhôṭpôṭ} — \textquotedblleft{}in a jiffy; quickly\textquotedblright{}.

\begin{quote}
\bn{মেয়েটা} \bn{ঝটপট} \bn{জামা} \bn{পরল।}
\end{quote}

\begin{quote}
\emph{Meyeṭā jhôṭpôṭ jāmā pôrlô.}
\end{quote}

\begin{quote}
The girl put on her dress in a jiffy.
\end{quote}

\begin{cousinweb}
\textbf{\bn{ঝটপট}} (\emph{jhôṭpôṭ}) is an expressive doubled word built on \bn{ঝট}, \textquotedblleft{}quickly\textquotedblright{}; CDIAL lists Bengali \emph{jhaṭ} and the doubled \emph{jhaṭ-phaṭ}, both \textquotedblleft{}quickly\textquotedblright{}, under a reconstructed \emph{jhaṭ}, \textquotedblleft{}sudden movement\textquotedblright{} (CDIAL 5327).
\end{cousinweb}

\begin{grammarlens}[title={What you've built}]
One more word for saying what happens if.
\end{grammarlens}

\subsection*{Your turn}
\begin{itemize}
\raggedright
  \item \emph{Say it:} \emph{jhôṭpôṭ}
  \item \emph{Say it:} \emph{jhôṭpôṭ}, once more
  \item \emph{Say it:} \emph{Meyeṭā jhôṭpôṭ jāmā pôrlô.} — The girl put on her dress in a jiffy.
  \item \emph{Recall:} say the Bengali for to draw near, then the Bengali for a scramble, then say \emph{jhôṭpôṭ} again
\end{itemize}

\begin{tcolorbox}[breakable,colback=teal!4,colframe=teal!35!black,title={Before you move on}]
What does \bn{ঝটপট} mean? (\textquotedblleft{}In a jiffy; quickly\textquotedblright{}.) Say it once more.
\end{tcolorbox}
